% !PRINTED 302
\weberpage{302}
If one applies the first of these formulae $\alpha_1$ times, the second $\alpha_2$ times
in succession, and so on, it follows that:

\[
\begin{aligned}
\varphi_1(z_1+\alpha_1, z_2, \dots) &\equiv \varphi_1(z_1, z_2, \dots) + \alpha_1\alpha_{1,1} \\
\varphi_1(z_1, z_2+\alpha_2, \dots) &\equiv \varphi_1(z_1, z_2, \dots) + \alpha_2\alpha_{1,2} \\
&\cdots\cdots\cdots\cdots\cdots\cdots\cdots\cdots\cdots,
\end{aligned}
\]

and if in the first one changes $z_2$ into $z_2+\alpha_2$, and applies the
second, and continues in this way, one finally obtains:

\weq{4}{\begin{aligned}
\varphi_1(z_1+\alpha_1, z_2+\alpha_2, \dots) &\equiv \\
\varphi_1(z_1, z_2, \dots) + \alpha_1\alpha_{1,1} + \alpha_2\alpha_{1,2} + \cdots
\end{aligned}}

In this we now set $z_1, z_2, \cdots = 0$ and then write in place of
$\alpha_1, \alpha_2, \dots$ again $z_1, z_2, \dots$. Further, we denote
$\varphi_1(0, 0, \dots)$ by $\alpha_1$ and thus obtain from (4):

\weq{5}{\varphi_1(z_1, z_2, \dots) \equiv \alpha_1 + z_1\alpha_{1,1} + z_2\alpha_{1,2} + \cdots,}

that is, $\varphi_1$ must be a \gesperrt{linear function}.

In exactly the same way we proceed with all the congruences (2) and thus arrive at the following final result:

\gesperrt{I. All permutations of the Galois group $P$ of a primitive irreducible metacyclic equation of degree $p^k$}

\[
\begin{pmatrix} z_1, z_2, \dots, z_k \\ z_1', z_2', \dots, z_k' \end{pmatrix}
\]

\gesperrt{are of the form}

\weq{6}{\begin{aligned}
z_1' &\equiv \alpha_{1,1}z_1 + \alpha_{1,2}z_2 + \cdots + \alpha_{1,k}z_k + \alpha_1 \\
z_2' &\equiv \alpha_{2,1}z_1 + \alpha_{2,2}z_2 + \cdots + \alpha_{2,k}z_k + \alpha_2 \quad (\mathrm{mod}\ p). \\
&\cdots\cdots\cdots\cdots\cdots\cdots\cdots\cdots\cdots\cdots
\end{aligned}}

\[
z_k' \equiv \alpha_{k,1}z_1 + \alpha_{k,2}z_2 + \cdots + \alpha_{k,k}z_k + \alpha_k
\]

The system of congruences (6) represents a permutation always and only when the determinant

\[
\sum \pm \alpha_{1,1}\alpha_{2,2}\cdots\alpha_{k,k}
\]

is not divisible by $p$, because then and only then does there also conversely belong to every system of the $z'$ a determinate system of the $z$ (with respect to the modulus $p$).

The totality of all permutations (6) is in fact a group, which is called the \gesperrt{general linear congruence group}, and in it the metacyclic groups $P$ must be con-
%CONT>>
